% =============================================================================
% 2.3 Analyse des besoins et spécification
% =============================================================================
\section{Analyse des besoins et spécification}
\label{sec:ch2-besoins}

\subsection{Besoins fonctionnels}

L'analyse des besoins a été menée en collaboration avec des praticiens du
domaine des appels d'offres au sein de SITINFRA Cameroun. Elle a permis
d'identifier les besoins fonctionnels du tableau~\ref{tab:besoins-f}.

\begin{btptable}{Besoins fonctionnels}
\label{tab:besoins-f}
\begin{tabularx}{\linewidth}{@{}C{1.1cm} L{3.6cm} C{2.0cm} Y@{}}
\toprule
\rowcolor{btppale}
\btpth{Id.} & \btpth{Besoin} & \btpth{Priorité} & \btpth{Description} \\
\midrule
BF1  & Ingestion de documents     & Haute    & Dépôt de fichiers (PDF, DOCX, tableurs, scans). \\
BF2  & Indexation                 & Haute    & Indexation sémantique et lexicale des extraits. \\
BF3  & Recherche                  & Haute    & Recherche hybride, filtrée par organisation et type. \\
BF4  & Extraction d'exigences     & Haute    & Extraction structurée des exigences du DAO. \\
BF5  & Génération de sections     & Haute    & Rédaction assistée d'une section, avec citations. \\
BF6  & Validation humaine         & Haute    & Revue, correction ou relance avant intégration. \\
BF7  & Export                     & Moyenne  & Assemblage et export du dossier final. \\
BF8  & Graphe de connaissances    & Basse    & Liens entre entités (seconde phase). \\
BF9  & Cloisonnement              & Moyenne  & Isolation des données par organisation. \\
BF10 & Historique                 & Moyenne  & Versions des sections et trace des décisions. \\
\bottomrule
\end{tabularx}
\end{btptable}

\subsection{Besoins non fonctionnels}

\begin{btptable}{Besoins non fonctionnels}
\label{tab:besoins-nf}
\begin{tabularx}{\linewidth}{@{}C{1.3cm} L{2.8cm} Y@{}}
\toprule
\rowcolor{btppale}
\btpth{Id.} & \btpth{Besoin} & \btpth{Description} \\
\midrule
BNF1 & Performance     & Temps de réponse inférieur à cinq secondes pour une recherche. \\
BNF2 & Fiabilité       & Taux d'erreur opérationnelle inférieur à un pour cent. \\
BNF3 & Sécurité        & Authentification, cloisonnement strict des corpus. \\
BNF4 & Traçabilité     & Chaque génération est rattachée à des sources et à une version. \\
BNF5 & Évolutivité     & Architecture modulaire, ajout de briques sans rupture. \\
BNF6 & Maintenabilité  & Code structuré, documenté, couvert par des tests. \\
\bottomrule
\end{tabularx}
\end{btptable}

\subsection{Spécifications fonctionnelles}

\paragraph{Ingestion et indexation (BF1, BF2).}
Le dépôt d'un fichier via l'interface web déclenche une détection du type
(numérique ou numérisé), une extraction préservant titres et tableaux, un
découpage hiérarchique parent-enfant, puis une indexation conjointe dans
un magasin vectoriel et une base relationnelle. Le document n'est marqué
« prêt » que lorsque les deux magasins ont confirmé l'écriture.

\paragraph{Extraction d'exigences (BF4).}
Le dossier d'appel d'offres courant est analysé pour en extraire les
exigences --- catégorie, caractère obligatoire, pièce source --- stockées
en base relationnelle et liées au dossier en cours. Cette extraction est
contrainte par un schéma~: le modèle ne produit pas un texte, il produit
des enregistrements validables.

\paragraph{Génération de sections (BF5, BF6).}
Un plan est d'abord construit à partir du dossier entrant. Pour chaque
section, le système construit une requête, récupère un contexte, génère
une sortie structurée, la valide automatiquement, puis s'interrompt pour
la revue humaine. Aucune section n'est intégrée au livrable sans décision
explicite.

\subsection{Critères d'évaluation}

Les critères du tableau~\ref{tab:criteres-eval} seront utilisés au
chapitre~\ref{chap:resultats} pour valider le système. Ils dérivent des
hypothèses H1 à H5 formulées en introduction générale.

\begin{btptable}{Critères d'évaluation et cibles}
\label{tab:criteres-eval}
\begin{tabularx}{\linewidth}{@{}L{3.4cm} Y C{2.6cm}@{}}
\toprule
\rowcolor{btppale}
\btpth{Critère} & \btpth{Mesure} & \btpth{Cible} \\
\midrule
Rappel à 20        & Proportion d'extraits pertinents dans les vingt premiers. & $> 0{,}75$ \\
Fidélité           & Proportion d'affirmations fondées sur le contexte. & $> 0{,}85$ \\
Temps d'ingestion  & Document numérique de cinquante pages. & $< 60$\,s \\
Temps de génération & Une section, hors revue humaine. & $< 30$\,s \\
Taux de validation & Sections approuvées sans modification. & $> 70$\,\% \\
\bottomrule
\end{tabularx}
\end{btptable}
